\title{Syllabus for Computer Science I: COSC120}
\author{Dr. Jordan Hanson - Whittier College Dept. of Physics and Astronomy}
\date{Fall 2026}
\documentclass[10pt]{article}
\usepackage[a4paper, total={18cm, 27cm}]{geometry}
\usepackage{outlines}
\usepackage{hyperref}
\begin{document}
\maketitle

\begin{abstract}
Computer Science I introduces computational problem solving and computer programming using Python. The course begins with statements, expressions, variables, primitive data types, input/output, and the relationship between variables and objects. Branching, Boolean logic, iteration, functions, and modules provide the foundation for structured program design. Strings, lists, tuples, and dictionaries are then used to represent and manipulate collections of data. The course continues with object-oriented programming, classes, recursion, inheritance, file input/output, and exception handling. Along the way, students apply Python to numerical computation, searching and sorting, simulation, testing and debugging, and elementary algorithmic analysis. The semester concludes with an introduction to data science using NumPy, Pandas, exploratory analysis, and visualization. Programming exercises and a final project emphasize decomposition, abstraction, verification, readable code, and communication of computational results.
\end{abstract}

\noindent
\textit{\textbf{Pre-requisites}: None. Prior programming experience is not required.} \\
\textit{\textbf{Course credits, Liberal Arts Categorization}: 3 Credits, LLB2, QR2} \\
\textit{\textbf{Regular course hours}: Mondays and Wednesdays from 8:30--9:50 AM in SLC 416.} \\
\textit{\textbf{Instructor contact information}: Discord: 918particle, Email: jhanson2@whittier.edu, Office: SLC 222, Book appointments: \url{https://calendar.app.google/gYfm8uzRbZ7fW9Dc7}.} \\
\noindent
\textit{\textbf{Office hours}: Use booking link above to schedule office hours.} \\
\textit{\textbf{Attendance/Absence}: Regular attendance helps everyone learn well, especially in team-coding and other group activities.  The instructor will collect attendance information using warm-up exercises.  Three absences for medical, family, or other legitimate reasons are permitted.  Student athletes complete a form with the instructor with information on the athletic schedule.  Unexcused absences will result in a grading penalty, but these parameters will be set by the students at the beginning of the course.} \\
\textit{\textbf{Late work policy}: Late work is accepted if the student and instructor agree on a modified deadline, within reason.} \\
\textit{\textbf{Text}: \textit{Introduction to Python Programming}, by Udayan Das, Aubrey Lawson, Chris Mayfield, and Narges Norouzi (OpenStax). The text is free and open-access, available at \\ \url{https://openstax.org/books/introduction-python-programming/pages/1-introduction}.} \\
\textit{\textbf{Software}: Python3. Students will use the OpenStax code runner and/or a local Python environment. Jupyter notebooks may also be used for numerical, simulation, and data-science exercises.} \\
\textit{\textbf{Grading}: The course grade will be a weighted average of assignment scores, and the weights are listed in Tab. \ref{tab:grading}.}
\begin{table}[hb]
\centering
\begin{tabular}{| l | l | l |}
\hline
\textbf{Assignment} & \textbf{Weight} & \textbf{Date} \\ \hline
Homework and programming sets & 40\% & Weekly on Fridays \\ \hline
First Midterm Exam & 20\% & Wednesday, September 30th \\ \hline
Second Midterm Exam & 20\% & Wednesday, November 4th \\ \hline
Final programming-project presentation & 20\% & Thursday, December 10th \\ \hline
\end{tabular}
\caption{\label{tab:grading} Major assignments and grade weights. Programming sets include short written exercises, code submissions, testing/debugging activities, and selected textbook problems.}
\end{table}
\noindent
\textit{\textbf{Grade Settings}: $<60\%$ = F, $\geq 60\%,< 70\%$ = D, $\geq 70\%,<80\%$ = C, $\geq 80\%,<90\%$ = B, $>90\%$ = A.  Pluses and minuses: 0-3\% minus, 3\%-6\% straight, 6\%-10\% plus (e.g. 79\% = C+, 91\% = A-)} \\
\textit{\textbf{Homework and Programming Sets}: Due weekly on Fridays. Assignments will combine textbook exercises, short programming problems, debugging, code-reading exercises, and larger applications. Students are expected to submit readable code with meaningful variable names and appropriate comments.} \\
\textit{\textbf{Final Programming Project}: Students will design, implement, test, and present a Python3 program that integrates several course concepts. Projects should demonstrate decomposition with functions and/or classes, use at least one nontrivial data structure, read or generate data, and communicate results clearly.} \\
\textit{\textbf{Student Accessibility Services (SAS) statement}: Whittier College is committed to making learning experiences as accessible as possible. If you experience physical or academic barriers due to a disability, you are encouraged to contact Student Accessibility Services (SAS) to discuss options. To learn more about academic accommodations, email sas@whittier.edu, call (562) 907-4825, or visit the SAS office on the ground floor of Wardman Library.} \\
\textit{\textbf{Academic Honesty Policy}: \url{https://www.whittier.edu/academics/academichonesty}} \\
\clearpage
\noindent
\textit{\textbf{Course Objectives}:}
\begin{itemize}
\item Develop Python3 fluency, including the use of data types, structures, conditional logic, loops, and algorithms.
\item Decompose programs into reusable modules, with reasonable parameters, return values, scope, and control flow.
\item Use strings, lists, tuples, and dictionaries to organize, search, transform, and analyze data.
\item Explain elementary algorithmic complexity and implement representative searching and sorting methods.
\item Apply numerical programming techniques to solve mathematical, statistical, and engineering problems.
\item Design programs using object-oriented programming, with classes and objects, encapsulation, and inheritance.
\item Read and write files, work with CSV data, and use exception handling to make programs more robust.
\item Test, debug, document, and communicate computational work clearly.
\end{itemize}
\textit{\textbf{Course Outline}:}
\begin{outline}[enumerate]
\1 Week 0 - \textbf{Reading: Chapter 1, Statements.}
\2 First meeting - Wednesday, August 19th, 2026.
\3 Course introduction and syllabus distribution
\3 Course motivation: coding is awesome!
\3 Python3 programs, statements, and the interpreter
\3 Input/output, variables, strings, numbers, comments, and error messages
\3 \textit{Programming exercise}
\4 Run Python3 programs using the OpenStax code runner
\4 Run Python3 programs using Google Colab
\1 Week 1 - \textbf{Reading: Chapters 2--3, Expressions and Objects.}
\2 First meeting - Monday, August 24th, 2026. \textit{Reading: Chapter 2.}
\3 Expressions, literals, operators, type conversion, and mixed data types
\3 Floating-point representation and floating-point error
\3 Integer division, remainder, and the \texttt{math} module
\3 \textit{Programming exercise, part I}
\4 Numerical expressions and unit-conversion programs
\4 Experiments with floating-point precision
\2 Second meeting - Wednesday, August 26th, 2026. \textit{Reading: Chapter 3.}
\3 Objects and references; strings revisited
\3 Lists and tuples as containers
\3 Formatted strings and formatted numerical output
\3 Binary, hexadecimal, and ASCII representations
\3 \textit{Programming exercise, part II}
\4 Decimal, binary, and hexadecimal conversion in Python
\4 String encoding and character-code experiments
\1 Week 2 - \textbf{Reading: Chapters 4--5, Decisions and Loops.}
\2 First meeting - Monday, August 31st, 2026. \textit{Reading: Chapter 4.}
\3 Boolean values and comparison operators
\3 \texttt{if}, \texttt{elif}, and \texttt{else}
\3 Boolean operations, precedence, chained decisions, and nested decisions
\3 \textit{Programming exercise, part I}
\4 Classification and decision-making programs
\4 Testing boundary cases and logical conditions
\2 Second meeting - Wednesday, September 2nd, 2026. \textit{Reading: Chapter 5.}
\3 \texttt{while} loops and \texttt{for} loops
\3 Nested loops, \texttt{break}, \texttt{continue}, and loop control
\3 Common iterative program patterns: counters, accumulators, sentinels, and validation
\3 \textit{Programming exercise, part II}
\4 Iterative numerical calculations
\4 Pattern generation and simple Monte Carlo experiments
\1 Week 3 - \textbf{Reading: Chapters 5--6, Loops and Functions.}
\2 Monday, September 7th, 2026. \textit{Labor Day: no instruction.}
\2 Meeting - Wednesday, September 9th, 2026. \textit{Reading: Review Chapter 5; Chapter 6.}
\3 Defining and calling functions
\3 Decomposition and abstraction using functions
\3 Parameters, arguments, return values, control flow, and variable scope
\3 Keyword arguments and reusable numerical functions
\3 Numerical programming: successive approximation and bisection methods
\3 \textit{Programming exercise}
\4 Refactor an iterative program into small reusable functions and test known inputs/outputs
\4 Implement square-root or root-finding by iteration; compare stopping criteria and floating-point limitations
\1 Week 4 - \textbf{Reading: Chapters 6--7, Functions and Modules.}
\2 First meeting - Monday, September 14th, 2026. \textit{Reading: Complete Chapter 6.}
\3 Function design, interfaces, documentation, and testing
\3 Scope, side effects, and debugging function-based programs
\3 \textit{Programming exercise, part I}
\4 Build a small multi-function numerical program
\4 Develop test cases before final submission
\2 Second meeting - Wednesday, September 16th, 2026. \textit{Reading: Chapter 7.}
\3 Module basics and importing names
\3 Top-level code and the \texttt{if \_\_name\_\_ == "\_\_main\_\_"} pattern
\3 Standard-library modules and the \texttt{help()} function
\3 \textit{Programming exercise, part II}
\4 Create and import a custom Python module
\4 Use selected standard-library modules for random numbers and mathematics
\1 Week 5 - \textbf{Reading: Chapters 8--9, Strings and Lists.}
\2 First meeting - Monday, September 21st, 2026. \textit{Reading: Chapter 8.}
\3 String operations, indexing, slicing, searching, and testing
\3 String formatting, splitting, and joining
\3 \textit{Programming exercise, part I}
\4 Text parsing and simple data-cleaning problems
\4 Palindrome and frequency-analysis examples
\2 Second meeting - Wednesday, September 23rd, 2026. \textit{Reading: Chapter 9.}
\3 Lists, mutability, iteration, sorting, and reversing
\3 Common list operations, nested lists, and list comprehensions
\3 Introduction to searching and sorting as algorithms
\3 \textit{Programming exercise, part II}
\4 Implement linear search and compare with Python3 list operations
\4 Implement insertion sort or selection sort and test on generated data
\1 Week 6 - \textbf{Reading: Review Chapters 1--9.}
\2 First meeting - Monday, September 28th, 2026.
\3 Review: statements, expressions, objects, decisions, loops, functions, modules, strings, and lists
\3 Strategies for testing and debugging
\3 Review of representative numerical and algorithmic problems
\3 \textit{Programming workshop}
\4 Catch-up, debugging, and midterm preparation
\2 Second meeting - Wednesday, September 30th, 2026.
\3 \textbf{First midterm exam}
\4 Covers Chapters 1--9 and associated programming exercises
\1 Week 7 - \textbf{Reading: Chapter 10, Dictionaries.}
\2 First meeting - Monday, October 5th, 2026. \textit{Reading: Sections 10.1--10.3.}
\3 Dictionary basics and key-value data
\3 Dictionary creation and common operations
\3 Choosing lists, tuples, or dictionaries for a problem
\3 \textit{Programming exercise, part I}
\4 Frequency tables and lookup problems
\4 Build a dictionary from structured input data
\2 Second meeting - Wednesday, October 7th, 2026. \textit{Reading: Sections 10.4--10.5.}
\3 Conditionals and loops with dictionaries
\3 Nested dictionaries and dictionary comprehensions
\3 Elementary algorithmic complexity: constant, logarithmic, linear, and quadratic growth
\3 Linear search and binary search; why sorted data changes the problem
\3 \textit{Programming exercise, part II}
\4 Compare linear search and binary search experimentally
\4 Measure operation counts as input size increases
\1 Week 8 - \textbf{Reading: Chapter 11, Classes.}
\2 Monday, October 12th, 2026. \textit{Indigenous Peoples' Day: no instruction.}
\2 Meeting - Wednesday, October 14th, 2026. \textit{Reading: Sections 11.1--11.5.}
\3 Object-oriented programming and abstract data types
\3 Classes, instances, attributes, and instance methods
\3 Abstraction, encapsulation, operator overloading, and class design
\3 Organizing classes in modules; interfaces and code reuse
\3 \textit{Programming exercise}
\4 Design a simple class to model a real-world entity; create and test multiple instances
\4 Extend the class with useful operators and validation; begin planning the final programming project
\1 Week 9 - \textbf{Reading: Chapter 12, Recursion.}
\2 First meeting - Monday, October 19th, 2026. \textit{Reading: Sections 12.1--12.3.}
\3 What is recursion? Base cases and recursive cases
\3 Recursive factorial, palindrome, and Fibonacci examples
\3 Recursion with strings and lists
\3 \textit{Programming exercise, part I}
\4 Trace recursive function calls by hand and in Python
\4 Compare recursive and iterative solutions
\2 Second meeting - Wednesday, October 21st, 2026. \textit{Reading: Sections 12.4--12.5.}
\3 Recursive problem solving and divide-and-conquer thinking
\3 Repeated subproblems, and avoiding repetition of work
\3 Merge sort as a recursive divide-and-conquer example
\3 \textit{Programming exercise, part II}
\4 Implement the Fibonacci sequence, or a similar recursive computation
\4 Implement or trace merge sort and compare with quadratic sorting methods
\1 Week 10 - \textbf{Reading: Chapter 13, Inheritance.}
\2 First meeting - Monday, October 26th, 2026. \textit{Reading: Sections 13.1--13.3.}
\3 Inheritance basics and relationships
\3 Attribute access and inherited methods
\3 Method overriding and extension
\3 \textit{Programming exercise, part I}
\4 Build a base class and one or more derived classes
\4 Test inherited and overridden behavior
\2 Second meeting - Wednesday, October 28th, 2026. \textit{Reading: Sections 13.4--13.5; begin Chapter 14.}
\3 Hierarchical inheritance and multiple inheritance
\3 Design tradeoffs, and ``mixin'' classes
\3 Introduction to persistent data and files
\3 \textit{Programming exercise, part II}
\4 Extend the Week 8 class design using inheritance
\4 Save simple program output to a text file
\1 Week 11 - \textbf{Reading: Chapter 14, Files and Exceptions; review Chapters 8--13.}
\2 First meeting - Monday, November 2nd, 2026. \textit{Reading: Sections 14.1--14.3.}
\3 Reading from files and writing to files
\3 Paths and files in different locations
\3 CSV files and structured tabular data
\3 \textit{Programming exercise}
\4 Read, process, and summarize a CSV data file
\4 Write processed results to a new file
\2 Second meeting - Wednesday, November 4th, 2026.
\3 \textbf{Second midterm exam}
\4 Covers Chapters 8--13 and associated programming exercises
\1 Week 12 - \textbf{Reading: Complete Chapter 14; applied computation.}
\2 First meeting - Monday, November 9th, 2026. \textit{Reading: Sections 14.4--14.5.}
\3 Exceptions and exception handling
\3 \texttt{try}, \texttt{except}, \texttt{else}, and \texttt{finally}
\3 Raising exceptions and validating program state
\3 \textit{Programming exercise, part I}
\4 Add robust input and file handling to an earlier program
\4 Design useful error messages and failure tests
\2 Second meeting - Wednesday, November 11th, 2026.
\3 Computation as a problem-solving tool
\3 Computer models of real-world phenomena
\3 Random processes, random walks, and simulation
\3 \textit{Programming exercise, part II}
\4 Random-walk simulation, thermodynamics calculations
\4 Collect simulation results for later visualization
\1 Week 13 - \textbf{Reading: Chapter 15, Data Science.}
\2 First meeting - Monday, November 16th, 2026. \textit{Reading: Sections 15.1--15.3.}
\3 Introduction to the data-science life cycle
\3 NumPy arrays and numerical data
\3 Pandas data structures and tabular data
\3 \textit{Programming exercise, part I}
\4 Load simulation or CSV data into NumPy/Pandas
\4 Compute descriptive summaries and selected transformations
\2 Second meeting - Wednesday, November 18th, 2026. \textit{Reading: Sections 15.4--15.5.}
\3 Exploratory data analysis
\3 Data filtering, indexing, and visualization
\3 Presenting computational results graphically
\3 \textit{Programming exercise, part II}
\4 Visualize simulation or tabular data with Matplotlib
\4 Interpret graphs and identify limitations of the analysis
\4 Final project development and proposal check
\1 Week 14 - \textbf{Fall/Thanksgiving Break:} November 23rd--27th, 2026. No instruction.
\1 Week 15 - \textbf{Reading: Review Chapter 15 and cumulative course concepts.}
\2 First meeting - Monday, November 30th, 2026.
\3 Integrating functions, data structures, classes, files, and visualization
\3 Final-project testing, debugging, and code review
\3 Next steps: courses offered at Whittier College, Life Lab 2 content
\3 \textit{Final project workshop}
\2 Second meeting - Wednesday, December 2nd, 2026.
\3 Course synthesis and cumulative review
\3 Student project demonstrations and peer feedback
\3 Professional presentation of computational results
\3 \textit{Final project workshop and presentation preparation}
\1 Week 16 - \textbf{Final instructional meeting.}
\2 Meeting - Monday, December 7th, 2026.
\3 Final project demonstrations and peer feedback
\3 Final project testing, presentation preparation, and cumulative course wrap-up
\1 Final Examination Period - \textbf{December 9th--12th, 2026.}
\2 Final programming-project presentations will on Thursday, December 10th in two sessions.
\end{outline}
\end{document}
